\section{Eventual sharp rank over arbitrary algebraically closed fields}
\label{sec:spreading}

This section transfers the finite-field rank theorem to every mixed pair by spreading the pair and using an open Frobenius-rank condition.

We use standard limit and cohomology results in their scheme-theoretic forms. Schemes of finite presentation, their morphisms and line bundles descend along filtered limits; smoothness, properness and relative ampleness hold after shrinking a finite-type model \citep[Tags~01ZM, 0B8W, 0C0C, 081F and 09MT]{Stacks}. For a proper finitely presented family and a finitely presented sheaf flat over the base, derived cohomology is perfect and commutes with arbitrary base change \citep[Tag~0B91]{Stacks}. Euler characteristic is locally constant \citep[Tag~0B9T]{Stacks}. Relative Serre vanishing supplies one threshold for all sufficiently large powers, and the resulting degree-zero direct image is locally free with arbitrary base change \citep[Tags~02O1 and 0D4E]{Stacks}.

\begin{lemma}\label{lem:stable-rank-open}
Let $\mathcal H$ be a vector bundle of rank $g$ on a scheme $T$ of characteristic two, with a map $\Phi:F_T^*\mathcal H\to\mathcal H$. The locus where the stable Frobenius rank on a geometric fibre is at least $q$ is open. It is the rank-at-least-$q$ locus of
\[
 \Phi^{[g]}:(F_T^g)^*\mathcal H\longrightarrow\mathcal H.
\]
In particular, $2$-rank at least $q$ is an open condition in a smooth proper family of geometrically connected curves of fixed genus over a reduced base.
\end{lemma}
\begin{proof}
For a semilinear endomorphism $F$ of a $g$-dimensional vector space over a perfect field, the subspaces $F^j(V)$ decrease. Once two consecutive subspaces agree, all subsequent ones agree. There can be at most $g$ strict dimension drops, so the stable image is $F^g(V)$. The asserted locus is therefore defined by nonvanishing of minors of the displayed bundle map.

For the curve assertion, put $\mathcal H=R^1f_*\OO_{\mathcal C}$. Its local freeness and base change follow from perfect cohomology and the constant fibre dimensions $h^0=1$, $h^1=g$, $h^i=0$ for $i>1$. Here is the elementary passage from these facts to local freeness. Represent the perfect complex locally by a bounded complex of finite free modules. Starting at its lowest term, the constant cohomology dimensions determine the constant ranks of the differentials successively. Over a reduced base, a matrix of constant fibre rank has locally an invertible maximal minor and zero complementary block: its larger minors vanish at every point and therefore vanish as functions. Kernels, images and cokernels are locally free, and the complex splits into its cohomology and contractible summands locally. This also proves base change. Relative Frobenius gives the indicated map on $\mathcal H$, and stable coherent Frobenius rank is the curve's $2$-rank.
\end{proof}

\begin{proof}[Proof of Theorem~\ref{thm:eventual-all-mixed}]
Choose a finitely generated field $K\subset k$ over $\FF_2$ over which $(A,M)$ is defined. After spreading and shrinking, there is an integral finite-type scheme $S/\FF_2$ with function field $K$, an abelian scheme $\pi:\mathcal A\to S$ of relative dimension two, and a relatively ample line bundle $\mathcal M$. Its geometric generic fibre, extended to $k$, is the original pair. The integer $\chi(\mathcal M_s)$ is the constant $D$.

We first arrange that every geometric fibre is mixed. By perfect cohomology, followed by splitting the complex at the generic point and shrinking, $\mathcal E=R^1\pi_*\OO_{\mathcal A}$ is a rank-two vector bundle commuting with base change. Frobenius gives
\[
 \Phi:F_S^*\mathcal E\longrightarrow\mathcal E.
\]
Its second iterate has rank one at the geometric generic point, by the preceding lemma and $f_A=1$. All its two-by-two minors vanish on integral $S$. Some entry is nonzero generically; restrict to a nonempty trivializing open where that entry is invertible. The second iterate now has rank one at every geometric point. Every fibre therefore has $2$-rank one.

Choose a closed point $s$ of this fixed nonempty finite-type scheme. Its residue field is finite, and the actual pair $(A_s,M_s)$ is defined over that field. Let $a_s=a(M_s)\ge0$. If $a_s>0$, Theorem~\ref{thm:finite-field} supplies an integer $e_s$ such that, for every $e\ge e_s$, the geometric system $|M_s^{3e}|$ has a smooth curve of defect two. If $a_s=0$, Theorem~\ref{thm:p10exact}(1) supplies a smooth curve of defect at most two for every $e\ge1$; take $e_s=1$. In both cases, for every $e\ge e_s$ there is a smooth curve $C_s$ satisfying
\begin{equation}\label{eq:special-rank-bound}
                    f_{C_s}\ge 9De^2-1.
\end{equation}
This argument does not require constancy of the ordinary index under specialization.

Relative Serre vanishing on the fixed Noetherian base gives an integer $e_1$ such that, for all $e\ge e_1$, the higher direct images of $\mathcal M^{3e}$ vanish and
\[
                  \mathcal E_e=\pi_*\mathcal M^{3e}
\]
is locally free and commutes with arbitrary base change. Enlarge $e_1$ so that these powers are relatively generated. This is one threshold on the fixed quasi-compact base, and it does not require deleting the chosen point $s$.

Let $P_e$ be the projective bundle of one-dimensional spaces of sections of $\mathcal E_e$. It is integral, and each geometric fibre is the full complete linear system of $M_s^{3e}$. Its tautological section defines a universal divisor. On each smooth integral abelian fibre a nonzero section is a nonzerodivisor; hence this is a relative effective Cartier divisor, flat over $P_e$ \citep[Tag~062Y]{Stacks}. Let $V_e\subset P_e$ be its smooth-fibre locus. This locus is open because the nonsmooth locus in the proper universal divisor has closed image. Its fibres are geometrically connected: on a geometric abelian fibre $A_t$ over an algebraically closed field $\Omega$, put $N=\mathcal M_t^{3e}$. Lemma~\ref{lem:cohomology}, or the sequence
\[
 0\longrightarrow N^{-1}\longrightarrow\OO_{A_t}
       \longrightarrow\OO_C\longrightarrow0,
\]
together with $H^0(N^{-1})=H^1(N^{-1})=0$, gives $H^0(\OO_C)=\Omega$. The genus is $g_e=9De^2+1$ by adjunction and Riemann--Roch.

Fix any $e\ge\max(e_s,e_1)$. The curve in~\eqref{eq:special-rank-bound} gives a geometric point of $V_e$ where $2$-rank is at least $g_e-2$. Lemma~\ref{lem:stable-rank-open} shows that this inequality defines a nonempty open $W_e\subset V_e$. Because $P_e$ is integral, $W_e$ contains its generic point and has a nonempty generic fibre over $S$. Extending that geometric generic fibre to $k$ gives a nonempty open in the original complete system $|M^{3e}|$ with
\[
                             f_C\ge g_e-2.
\]
The positive-index upper bound in Theorem~\ref{thm:p10bounds} gives the reverse inequality, since $a(M)>0$. Thus $f_C=g_e-2$ generically, and~\eqref{eq:dimensions} gives $r(C)=9De^2-2$.

The same point $s$ and the same two integers $e_s,e_1$ serve every $e\ge e_0=\max(e_s,e_1)$. Only the projective bundle and its open subset depend on $e$. There is no intersection of infinitely many open subsets of $S$. Since the original pair is the geometric generic fibre of its own spread, the conclusion holds for every pair in the theorem.
\end{proof}

The proof is existential for an arbitrary pair. For a pair already defined over a finite field, Theorem~\ref{thm:finite-field} retains its explicit sufficient bound. An index-zero ample bundle on the special fibre need not extend to the generic fibre; the proof transfers a nonvanishing Frobenius minor instead. Thus eventual sharp rank does not imply existence of an index-zero ample class on the original surface.
